\subsection{弗雷德霍姆映射的扰动}

定理 3. --- 设 E 是局部凸空间，F 是分离局部凸空间，u 是从 E 到 F 中的弗雷德霍姆映射，h 是从 E 到 F 中的紧线性映射。则 \(u + h\) 是弗雷德霍姆算子，且有 \(\operatorname{ind}(u + h) = \operatorname{ind}(u)\) 。

我们将分三步证明该定理。

\begin{enumerate}
\def\labelenumi{\Alph{enumi})}
\item
  假设 E 和 F 是巴拿赫空间。III, p. 42 的命题 2 表明线性映射 u 是具有闭像且核与余核均有限维的严格态射。由于 E 和 F 是巴拿赫空间，III, p. 73 的定理 1 与定理 2 蕴含：对每个 \(t \in [0,1]\) ，线性映射 \(u_t = u + th\) 具有这些同样的性质，从而是弗雷德霍姆算子（III, p. 42, 命题 2）。映射 \(t \mapsto u_t\) （从 \([0,1]\) 到从 E 到 F 的弗雷德霍姆映射所成的集 \(\mathcal{F}(\mathrm{E};\mathrm{F})\) 中）是连续的。由 III, p. 58 的定理 1，映射 \(t \mapsto \operatorname{ind}(u_t)\) 是局部常数的。由于区间 \([0,1]\) 连通，这个映射是常数。于是我们有 \(\operatorname{ind}(u) = \operatorname{ind}(u + h)\) 。
\item
  我们假设 E = F，\(u = 1_{E}\) 。由 III, p. 5 的命题 5，存在一个巴拿赫空间 G、一个连续线性映射 \(p: E \to G\) 和一个紧线性映射 \(q: G \to E\)，使得 \(h = q \circ p\) 。G 的自同态 \(p \circ q\) 是紧的。由 A)，\(1_{G} + p \circ q\) 是 G 的指标为零的弗雷德霍姆自同态。但此时 \(1_{E} + h = 1_{E} + q \circ p\) 是 E 的指标为零的弗雷德霍姆自同态（III, p. 49, 命题 10, f)）。
\item
  一般情形。设 \(v\) 是 \(u\) 的一个拟逆。自同态 \(u \circ v\) 与 \((u + h) \circ v\) （\(F\) 的）与 \(1_{F}\) 都只差一个紧线性映射。由 B)，它们是 \(F\) 的指标为零的弗雷德霍姆自同态。由此得出 \(u + h\) 是弗雷德霍姆自同态（III, p. 40, 第 2 目）且与 \(u\) 有相同的指标（III, p. 43, 第 3 目, 公式 (2)）。
\end{enumerate}

推论 1. --- 设 E 和 F 是分离局部凸空间，且 \(u \in \mathcal{L}(\mathrm{E};\mathrm{F})\) 。为使 u 是弗雷德霍姆算子，必须且只需存在一个映射 \(v \in \mathcal{L}(\mathrm{F};\mathrm{E})\) ，使得线性映射 \(1_{\mathrm{E}} - v \circ u\) 与 \(1_{\mathrm{F}} - u \circ v\) 都是紧的。

由于有限秩的连续线性映射是紧的，所述条件是必要的。

设 \(v\) 是 \(\mathcal{L}(\mathrm{F};\mathrm{E})\) 中使得线性映射 \(1_{\mathrm{E}} - v \circ u\) 与 \(1_{\mathrm{F}} - u \circ v\) 为紧的的元素。由定理 3，\(v \circ u\) 与 \(u \circ v\) 是弗雷德霍姆映射。设 \(p\) 与 \(q\) 分别是 \(v \circ u\) 与 \(u \circ v\) 的拟逆。令 \(w = p \circ v\) ，\(w' = v \circ q\) 。采用 III, p. 40 第 2 目的记号，我们有 \(w \circ u \equiv 1_{\mathrm{E}}\) 与 \(u \circ w' \equiv 1_{\mathrm{F}}\) ，由此 \(w \equiv w \circ u \circ w' \equiv w'\) 。于是得出 \(w\) 是 \(u\) 的一个拟逆。

推论 2. --- 设 E 和 F 是分离局部凸空间，且 \(u \in \mathcal{L}(\mathrm{E};\mathrm{F})\) 。下列条件等价：

\begin{enumerate}
\def\labelenumi{(\roman{enumi})}
\item
  映射 u 是指标为零的弗雷德霍姆映射；
\item
  存在一个从 E 到 F 上的同构 v，使得 u - v 有限秩；
\item
  存在同构 \(v\) （从 \(E\) 到 \(F\) 上的），使得 \(u - v\) 是紧的。
\item
  \(\Longrightarrow\) (ii)：假设 u 是指标为零的弗雷德霍姆算子。存在拓扑直和分解 \(E = E_{1} \oplus E_{2}\) 与 \(F = F_{1} \oplus F_{2}\) （其中 \(E_{2}\) 与 \(F_{2}\) 有限维），使得 u 在 \(E_{2}\) 上为零，且通过限制定义一个同构 \(u_{1}\) （从 \(E_{1}\) 到 \(F_{1}\) 上）（III, p. 42, 命题 2）。若 \(\text{ind}(u) = 0\) ，则 \(E_{2}\) 与 \(F_{2}\) 的维数相等，且存在从 E 到 F 上的一个同构 v，它是 \(u_{1}\) 的扩张。u - v 的核包含 \(E_{1}\) ，故 u - v 有限秩。
\end{enumerate}

由于从 E 到 F 中的每个有限秩连续线性映射都是紧的，我们有 (ii) \(\Longrightarrow\) (iii)；而由定理 3 有 (iii) \(\Longrightarrow\) (i)。

